% Mechanical relocation generated from the preservation ledger.
% Existing prose is included byte-for-byte from preserved blocks.

\section{Discussion}

This software framework combines timetable information and operational counts to infer OD demand and quantify uncertainty,
while using a differentiable assignment model compatible with ML, MAP, and
gradient-based Bayesian inference. ML and MAP provide computationally convenient
point estimates and local curvature diagnostics, whereas Bayesian VI additionally
provides an approximate distribution that represents parameter uncertainty.

Planned extensions of the software include incorporating capacity effects in a
differentiable manner, estimating additional measurement parameters
(e.g., $\rho$ and $r$), supporting weighted measurement aggregations,
and adding opt-out alternatives.


\section{Conclusions and Recommended Use}

The central difficulty is not only computational scale but intrinsic
underdetermination: boarding and alighting counts observe events along journeys
rather than complete passenger trajectories. A credible solution must combine
a transparent forward model, explicit structural zeros, a scientifically
defensible demand restriction, regularization, and validation at observable
counts.

For the full network, the recommended primary workflow is fixed-routing,
conditional-gravity MAP estimation. Fixed routing makes the response additive;
gravity replaces one parameter per OD-time cell with a small set of shared
behavioral and production coefficients; MAP stabilizes directions weakly
informed by counts. ML remains the principal prior-sensitivity reference, and
variational Bayesian inference is appropriate when posterior uncertainty is
required and its approximation can be validated.

This recommendation is conditional, not universal. Route-level IPF and
entropic transport remain useful audits and initializers, while cell-level MAP
is valuable on small examples and for justified local refinement. Final
acceptance requires reproducible optimization, adequate grouped residuals,
reasonable curvature and prior sensitivity, grouped holdout, and agreement
with detailed assignment under matching assumptions.

The next scientific extensions are evidence-led model relaxations, explicit
measurement-design analysis, and capacity- or crowding-aware routing where the
data require it. Until then, results should be reported as estimates under the
stated timetable, routing, production, attractiveness, likelihood, and prior
assumptions, with weakly informed quantities identified rather than hidden.
